% ---
% RESUMOS
% ---

% RESUMO em português
\setlength{\absparsep}{18pt} % ajusta o espaçamento dos parágrafos do resumo
\begin{resumo}
 Este trabalho compara, em simulação, três formas de pilotar um veículo de
 corrida autônomo em escala 1:10 da plataforma F1TENTH: controladores
 clássicos, políticas obtidas por aprendizado de máquina e pilotos humanos. Os
 pilotos clássicos são um seguidor de trajetória pure pursuit, que usa pose
 global e a linha de corrida e também gera as demonstrações de treino, e dois
 controladores reativos que usam apenas o LiDAR, o Disparity Extender e um
 seguidor de gap por descontinuidade. Os pilotos de aprendizado de máquina são
 uma rede de imitação convolucional unidimensional (TinyLidarNet), que dispensa
 pose global, e um resíduo de aprendizado por reforço (SAC) aplicado sobre ela.
 Os pilotos humanos usam um joystick, com os mesmos limites de ação e o mesmo
 atraso de esterço das políticas, em duas vistas: a pista inteira vista de cima
 e apenas os pontos do LiDAR. As políticas aprendidas foram aprovadas por
 critérios de aceite congelados antes da medição. O comparativo usa um protocolo
 comum na pista Spielberg, com randomização de atrito, de ruído do sensor e de
 largada, e mede conclusão de volta, tempo de volta, colisões, taxa de esterço,
 desvio da linha de corrida e perfil de velocidade, num estudo de caso
 descritivo, sem teste de hipótese. Nos portões de aceite, o pure pursuit
 completa 50 de 50 voltas em 59,000 s, sem randomização; a rede de imitação, 50
 de 50 em 60,050 s; e o resíduo, 91 de 92 em 52,525 s, 11,0\% mais rápido que o
 especialista. Sob o protocolo comum, o resíduo faz 52,55 s com 50 de 50
 voltas, o pure pursuit 59,00 s com 44 de 45 e a rede de imitação 60,05 s com
 50 de 50; o Disparity Extender completa 50 de 50 voltas em 145,55 s, com a
 maior taxa de esterço; e o seguidor de gap não completa volta em nenhuma das 48
 combinações de parâmetros varridas nem na avaliação, com 0 de 5 voltas.
 \pendentehumano{uma frase com o resultado dos pilotos humanos em cada vista}.
 Conclui-se que, com o mesmo sensor em execução, as políticas aprendidas
 alcançaram o desempenho do controlador que usa o mapa e os controladores
 reativos não o alcançaram, com a ressalva de que a informação do mapa entrou nas
 políticas aprendidas pelas demonstrações de treino.

 \textbf{Palavras-chave}: F1TENTH. Piloto humano. Controle reativo.
 Aprendizado por imitação. Aprendizado por reforço. Simulação.
\end{resumo}

% ABSTRACT in english
\begin{resumo}[Abstract]
 \begin{otherlanguage*}{english}
   This work compares, in simulation, three ways of driving a 1:10 scale
   autonomous racing vehicle of the F1TENTH platform: classical controllers,
   policies obtained by machine learning, and human drivers. The classical
   drivers are a pure pursuit path follower, which uses the global pose and
   the racing line and also generates the training demonstrations, and two
   reactive controllers that use only the LiDAR, the Disparity Extender and a
   discontinuity-based gap follower. The machine learning drivers are a
   one-dimensional convolutional imitation network (TinyLidarNet), which
   does not require a global pose, and a residual reinforcement learning
   policy (SAC) applied on top of it. The human drivers use a joystick, with
   the same action limits and the same steering delay as the policies, in
   two views: the whole track seen from above and only the LiDAR points. The
   learned policies were approved by acceptance criteria frozen before
   measurement. The comparison uses a common protocol on the Spielberg
   track, with randomized friction, sensor noise and start pose, and
   measures lap completion, lap time, collisions, steering rate, deviation
   from the racing line and speed profile, in a descriptive case study,
   without hypothesis testing. At the acceptance gates, pure pursuit
   completes 50 of 50 laps in 59.000 s, without randomization; the imitation
   network, 50 of 50 in 60.050 s; and the residual policy, 91 of 92 in
   52.525 s, 11.0\% faster than the expert. Under the common protocol, the
   residual policy runs 52.55 s with 50 of 50 laps, pure pursuit 59.00 s with
   44 of 45, and the imitation network 60.05 s with 50 of 50; the Disparity
   Extender completes 50 of 50 laps in 145.55 s, with the highest steering
   rate; and the gap follower completes no lap in any of the 48 swept
   parameter combinations nor in the evaluation, with 0 of 5 laps.
   \pendentehumano{uma frase com o resultado dos pilotos humanos em cada vista}.
   It is concluded that, with the same sensor at run time, the learned
   policies reached the performance of the controller that uses the map,
   while the reactive controllers did not, with the caveat that the map
   information entered the learned policies through the training
   demonstrations.

   \vspace{\onelineskip}

   \noindent
   \textbf{Keywords}: F1TENTH. Human driver. Reactive control. Imitation
   learning. Reinforcement learning. Simulation.
 \end{otherlanguage*}
\end{resumo}
